\subsection{The \texorpdfstring{$41$}{41} integrity index}

The integer $41$ is not an isolated numerical curiosity. On the retained branch it is best read as a weighted integrity index rather than as a topological Euler-type object.
\begin{definition}[The \texorpdfstring{$41$}{41} integrity index]
\label{def:integrity-index-41}
Let
\[
\mathcal I_{41}:=10\,\bone=\Omega_{\mathrm{occ}}\gamma+N_\Phi.
\]
On the canonical branch, $\Omega_{\mathrm{occ}}=\Oadm=48$ and $N_\Phi=1$, so
\[
\mathcal I_{41}=40+1=41.
\]
\end{definition}

Then the hard arithmetic chain reads
\begin{equation}
\label{eq:41-chain-main}
10\,\bone
=
\Oadm\gamma+1
=
41.
\end{equation}
\begin{figure}[t]
\centering
\footnotesize
\begin{tikzpicture}[
    >=Latex,
    chain/.style={draw=black!55, rounded corners, align=center, text width=2.6cm, minimum height=1.0cm, inner sep=4pt},
    optchain/.style={draw=purple!70!black, dashed, rounded corners, align=center, text width=4.2cm, minimum height=0.95cm, inner sep=4pt},
    eq/.style={font=\Large},
    note/.style={font=\scriptsize\itshape, text=black!55, align=center}
]
\node[chain, fill=green!7, draw=green!50!black] (a) at (0,0) {$10\,\bone$\\\scriptsize abelian index};
\node[eq] at (2.15,0) {$=$};
\node[chain, fill=orange!9, draw=orange!75!black] (b) at (4.3,0) {$\Oadm\gamma+1$\\\scriptsize occupancy $+$ Higgs};
\node[eq] at (6.45,0) {$=$};
\node[chain, fill=gray!8, draw=black!65] (c) at (8.6,0) {$41$\\\scriptsize integrity marker};
\draw[dashed, orange!60!black, ->, thick] (b.south) -- ++(0,-0.35);
\node[optchain, below=0.5cm of b] (d) {$48\,I_1^{\mathrm{EW}}=41$\\ on the EW normalization\\[1pt]$32\,k_{\mathrm{EW}}=41$\\ only in the appendix-level\\E$_8$ reading};
\node[note, below=0.15cm of d] {optional extension, not part of the hard chain};
\end{tikzpicture}
\caption{The $41$-chain strip. The hard arithmetic chain is $10\,\bone=\Oadm\gamma+1=41$. Electroweak normalization and E$_8$ block-language extensions are shown only as optional downstream readings.}
\label{fig:41-chain-strip}
\end{figure}
The extra $+1$ is precisely the arithmetic contribution of one light seam-even Higgs doublet. This makes $\mathcal I_{41}$ a hypercharge-driven occupancy integrity index:
\[
\mathcal I_{41}=40+1,
\]
with $40=\Oadm\gamma$ the fermionic weighted occupancy and $1=N_\Phi$ the light Higgs contribution. The canonical branch is therefore locked arithmetically at $40\mapsto 41$, while the actual dynamical Higgs-selection theorem still comes from the bosonic / Callias closure route.

\begin{remark}[Optional block-trace extension]
If one introduces the electroweak block normalization
\[
I_1^{\mathrm{EW}}:=\frac{5}{24}\,\bone,
\]
then
\[
48\,I_1^{\mathrm{EW}}=41.
\]
The further identity
\[
32\,k_{\mathrm{EW}}=41
\]
belongs only to the appendix-level E$_8$ reading and requires the block-index extrapolation.
\end{remark}

\begin{remark}[Integrity-index reading]
The canonical branch may be read as carrying an integrity marker: the hypercharge trace coefficient and the occupancy-weighted carrier norm plus one light Higgs doublet collapse onto the same integer $41$. This is not a cryptographic claim. It is an algebraic consistency signature showing that several independently defined carrier-side layers agree on one integer.
\end{remark}

\subsection{Single-seed compression and the Cabibbo inversion}

The seed-driven observation rows $(\beta_{\mathrm{rad}},\Omega_b,\lambda_C,\sin^2\theta_{13})$
are not four independent hits but four projections of $\phiz$. The Cabibbo angle admits an
explicit inversion back to the seed:
\begin{equation}
\label{eq:cabibbo-inversion}
\phiz=\frac{1-\sqrt{1-4\lambda_C^2}}{2}.
\end{equation}
\begin{remark}[Catalan expansion of the Cabibbo inversion]
Setting
\[
x:=\lambda_C^2,
\]
the inversion formula becomes
\[
\phiz=x\,C(x),
\qquad
C(x):=\frac{1-\sqrt{1-4x}}{2x}.
\]
Hence
\begin{equation}
\label{eq:catalan-cabibbo}
\phiz
=
\sum_{n\ge 0} C_n\,\lambda_C^{2n+2}
=
\lambda_C^2+\lambda_C^4+2\lambda_C^6+5\lambda_C^8+14\lambda_C^{10}+\cdots,
\end{equation}
where $C_n$ are the Catalan numbers. In the present paper the appropriate interpretation of this
identity is noncrossing combinatorics on the seed side; it is an exact algebraic strengthening of
the one-seed compression rather than a claim about planar Feynman graphs, a specific large-$N$ /
't Hooft resummation, or a holographic flavor limit.
\end{remark}

From this all other bridge outputs become functions of $\lambda_C$ alone:
\begin{equation}
\label{eq:lambda-c-master}
\begin{aligned}
\beta_{\mathrm{rad}}&=\frac{1-\sqrt{1-4\lambda_C^2}}{8\pi},\\
\Omega_b&=\frac{4\pi-1}{8\pi}\bigl(1-\sqrt{1-4\lambda_C^2}\bigr),\\
\sin^2\theta_{13}&=\frac{e^{-\gamma}}{2}\bigl(1-\sqrt{1-4\lambda_C^2}\bigr).
\end{aligned}
\end{equation}
\begin{remark}[Exact seed-budget identity]
The same one-seed map carries an exact budget identity:
\begin{equation}
\label{eq:seed-budget-identity}
\Omega_b+\beta_{\mathrm{rad}}=\phiz.
\end{equation}
Indeed,
\[
\Omega_b+\beta_{\mathrm{rad}}
=(4\pi-1)\beta_{\mathrm{rad}}+\beta_{\mathrm{rad}}
=4\pi\,\beta_{\mathrm{rad}}
=\phiz.
\]
This is best read as an exact seed-splitting identity inside the intrinsic observation layer, not
as an independent dynamical conservation law.
Because $\Omega_b$ and $\beta_{\mathrm{rad}}$ are dimensionless density fractions rather than angle measures, any $4\pi$ solid-angle reading of this identity remains interpretive rather than theorem-level.
\end{remark}

Thus on the canonical branch, Cabibbo alone pulls $\beta$, $\Omega_b$, and $\theta_{13}$. Conversely, $\theta_{13}$ pulls $\lambda_C$ and $\Omega_b$ through $\phiz=e^\gamma\sin^2\theta_{13}$. This is a much stronger correlation statement than the seed map alone.

\subsection{Carrier exclusion theorems}

The strongest argument that the carrier packet selects the physical world is not a further precision hit. It is that wrong discrete worlds fail brutally.

\begin{theorem}[Carrier exclusion]
\label{thm:carrier-exclusion}
Let $\alpha^{-1}(0;N_{\mathrm{fam}},N_\Phi,\nu^c)$ denote the positive self-consistent root of the carrier-form electromagnetic closure equation \eqref{eq:carrier-cfe} with the indicated discrete data. Then:
\begin{center}
\small
\sisetup{round-mode=none, group-minimum-digits=20}
\setlength{\tabcolsep}{2pt}
\renewcommand{\arraystretch}{0.96}
\begin{tabularx}{0.98\linewidth}{@{}>{\raggedright\arraybackslash}p{0.22\linewidth}>{\raggedright\arraybackslash}p{0.20\linewidth}>{\raggedleft\arraybackslash}p{0.14\linewidth}>{\raggedleft\arraybackslash}p{0.14\linewidth}Y@{}}
\toprule
Discrete world & Change from canonical & $\alpha^{-1}(0)$ & Residual vs CODATA & Reading \\
\midrule
\textbf{canonical $(3,1,\text{yes})$} & --- & $\mathbf{\num{137.0359992}}$ & $+\num{3.9e-8}$ & \makecell[l]{inside CODATA window\\($+1.87\sigma$)} \\
$N_{\mathrm{fam}}=2$ & one family fewer & $\num{156.094}$ & $+19.1$ & catastrophically excluded \\
$N_{\mathrm{fam}}=4$ & one family more & $\num{124.835}$ & $-12.2$ & catastrophically excluded \\
$N_\Phi=2$ & two light Higgs doublets & $\num{135.946}$ & $-1.09$ & catastrophically excluded \\
no $\nu^c$ & remove right-handed neutrino & $\num{137.0338}$ & $-\num{2.16e-3}$ & \makecell[l]{excluded\\(${\sim}\num{1.03e5}\sigma$)} \\
\bottomrule
\end{tabularx}
\end{center}
\end{theorem}

\begin{proof}[Proof sketch]
Hold fixed
\[
\cthree=\frac{1}{8\pi},
\qquad
\phibase=\frac{1}{6\pi},
\qquad
\gamma=\frac56,
\qquad
B=\frac32,
\]
and solve the positive self-consistent root of the same scalar closure equation on each discrete world:
\[
\alpha^3-2\cthree^3\alpha^2-8\cthree^6\,\bone(N_{\mathrm{fam}},N_\Phi)\ln\!\bigl(\phiseam(\alpha)^{-1}\bigr)=0,
\]
with
\[
\phiseam(\alpha)=\phibase+\deltatop e^{-2\alpha}(1-\deltatop e^{-2\alpha})^{-5/4},
\qquad
\deltatop=\Oadm\cthree^4,
\qquad
\delta_2=\frac54\,\deltatop^2,
\]
and
\[
\bone(N_{\mathrm{fam}},N_\Phi)=\frac43N_{\mathrm{fam}}+\frac1{10}N_\Phi.
\]
For the no-$\nu^c$ row, $\bone$ stays unchanged because $Y(\nu^c)=0$, while the occupancy changes from $\Oadm=48$ to $\Oadm=45$ because the one-family packet drops from $16$ to $15$ states. The same positive-branch solver and tolerance are used for every row, so the deviations are arithmetic consequences of changed discrete data rather than fits.
\end{proof}

\begin{remark}
Two families shift $\alpha^{-1}$ by $+19$, four families by $-12$, and two Higgs doublets by $-1.09$. Even the mildest alternative---removing $\nu^c$---shifts the value by $\num{2.16e-3}$, which is $\sim\num{1.03e5}$ times the CODATA uncertainty. This is not marginal discrimination; it is categorical exclusion. The canonical branch $(3,1,\text{yes})$ is the only discrete world that survives the electromagnetic closure test within the CODATA window.
\end{remark}

\begin{figure}[t]
\centering
\scriptsize
\begin{tikzpicture}[
    >=Latex,
    head/.style={draw=black!65, fill=gray!15, rounded corners, minimum height=0.8cm, align=center, font=\bfseries},
    ok/.style={draw=green!50!black, fill=green!8, rounded corners, minimum height=0.8cm, align=left},
    bad/.style={draw=red!70!black, fill=red!6, rounded corners, minimum height=0.8cm, align=left}
]
\node[head, minimum width=3.5cm] at (0,0) {discrete world};
\node[head, minimum width=2.8cm] at (4.15,0) {Residual vs CODATA};
\node[head, minimum width=4.6cm] at (8.2,0) {reading};

\node[ok, anchor=west, text width=3.25cm] at (-1.72,-1.0) {canonical $(3,1,\mathrm{yes})$};
\node[ok, minimum width=2.55cm] at (4.15,-1.0) {$+\num{3.9e-8}$};
\node[ok, anchor=west, text width=4.35cm] at (6.0,-1.0) {inside the CODATA window};

\node[bad, anchor=west, text width=3.25cm] at (-1.72,-1.95) {$N_{\mathrm{fam}}=2$};
\node[bad, minimum width=2.55cm] at (4.15,-1.95) {$+19.1$};
\node[bad, anchor=west, text width=4.35cm] at (6.0,-1.95) {catastrophically excluded};

\node[bad, anchor=west, text width=3.25cm] at (-1.72,-2.9) {$N_{\mathrm{fam}}=4$};
\node[bad, minimum width=2.55cm] at (4.15,-2.9) {$-12.2$};
\node[bad, anchor=west, text width=4.35cm] at (6.0,-2.9) {catastrophically excluded};

\node[bad, anchor=west, text width=3.25cm] at (-1.72,-3.85) {$N_\Phi=2$};
\node[bad, minimum width=2.55cm] at (4.15,-3.85) {$-1.09$};
\node[bad, anchor=west, text width=4.35cm] at (6.0,-3.85) {extra light doublet breaks closure};

\node[bad, anchor=west, text width=3.25cm] at (-1.72,-4.8) {no $\nu^c$};
\node[bad, minimum width=2.55cm] at (4.15,-4.8) {$-\num{2.16e-3}$};
\node[bad, anchor=west, text width=4.35cm] at (6.0,-4.8) {$\sim\num{1.03e5}\sigma$ exclusion};
\end{tikzpicture}
\caption{Exclusion matrix for wrong discrete worlds. The middle column records the residual versus the CODATA benchmark rather than a misleading shift relative to the canonical row. Once the discrete packet is changed, the electromagnetic closure leaves the observed window decisively.}
\label{fig:carrier-exclusion-matrix}
\end{figure}

\section{Closure architecture A: families, Higgs, and local spectral scale}

This first closure-architecture block is the shortest route from the hard carrier packet to a broader Standard Model plus gravity compiler. It has three moves: the local patch-deficit count gives the carrier-side family count, the topological family closure globalizes it into the family space $F$, the bosonic relative index selects the unique light Higgs doublet, and the local spectral-scale channel generates both $\Mpl$ and the electroweak vacuum from one common geometric response.

\subsection{Family monodromy and occupancy}

\begin{definition}[Family space]
Let $X_f$ denote the fundamental section of the seam complement. Define
\[
F:=H^1(X_f,\mathbb C).
\]
We use the distinguished cycle basis $(C_1,C_2,C_T)$ only as a concrete carrier-side basis of this family space; any later APS/Hodge realization is read as compatible analytic structure on the same topological family space.
\end{definition}
In this interpretation $F$ is neither an arbitrary flavor vector space nor a loose bookkeeping trick; it is a cycle space attached to the seam data. The actual bridge logic is shown in \Cref{fig:family-triangle}: the local patch-deficit count fixes $N_{\mathrm{fam}}=3$, the topological family step promotes that count to a family space, and the occupancy reinterpretation then feeds back into the defect channel.

\begin{figure}[t]
\centering
\scriptsize
\begin{tikzpicture}[
    >=Latex,
    hard/.style={draw=blue!60!black, fill=blue!6, rounded corners, align=center, text width=2.2cm, minimum height=1.0cm},
    closure/.style={draw=orange!75!black, fill=orange!8, rounded corners, align=center, text width=2.4cm, minimum height=1.0cm},
    bridge/.style={draw=green!50!black, fill=green!8, rounded corners, align=center, text width=2.25cm, minimum height=1.0cm}
]
\draw[draw=blue!60!black, fill=blue!4, rounded corners] (-5.5,-1.4) rectangle (-3.0,1.4);
\foreach \x/\y in {-5.1/1.0,-3.4/1.0,-5.1/-1.0,-3.4/-1.0} {
    \filldraw[blue!60!black] (\x,\y) circle (2.3pt);
    \node[blue!60!black, anchor=south] at (\x,\y) {\tiny $\frac34$};
}
\node[align=center] at (-4.25,0) {\textbf{local patch deficit}\\$F_{\mathrm{patch}}=4$ corners\\each contributes $\frac34$};

\node[hard] (nfam) at (-1.0,0) {$N_{\mathrm{fam}}$\\$=4\cdot\frac34=3$};

\coordinate (A) at (2.0,1.15);
\coordinate (B) at (0.7,-0.95);
\coordinate (C) at (3.3,-0.95);
\node[closure, text width=1.2cm, minimum height=0.8cm] at (A) {$C_1$};
\node[closure, text width=1.2cm, minimum height=0.8cm] at (B) {$C_2$};
\node[closure, text width=1.2cm, minimum height=0.8cm] at (C) {$C_T$};
\draw[->, thick, orange!75!black] ($(A)+(0.2,-0.3)$) .. controls (2.7,0.25) .. ($(C)+(-0.2,0.3)$);
\draw[->, thick, orange!75!black] ($(C)+(-0.38,-0.02)$) .. controls (2.0,-1.65) .. ($(B)+(0.38,-0.02)$);
\draw[->, thick, orange!75!black] ($(B)+(0.2,0.3)$) .. controls (1.3,0.2) .. ($(A)+(-0.2,-0.3)$);
\node[closure, text width=2.1cm, minimum height=0.9cm] at (2.0,-2.2) {topological family closure\\$F\cong\mathbb C^3$};

\node[closure] (occ) at (5.85,0.55) {$\Omega_{\mathrm{adm}}$\\$=3\times16=48$};
\node[bridge] (dtop) at (5.85,-1.15) {$\deltatop=\Omega_{\mathrm{adm}}\cthree^4$\\boundary-to-bulk bridge};

\draw[->, thick] (-3.0,0) -- (nfam);
\draw[->, thick] (nfam) -- (1.0,0.1);
\draw[->, thick] (nfam) -- (occ);
\draw[->, thick, orange!75!black] (2.0,-1.55) -- (2.0,-1.8);
\draw[->, thick, orange!75!black] (3.05,-1.9) -- (occ);
\draw[->, thick, green!50!black] (occ) -- (dtop);
\end{tikzpicture}
\caption{Carrier-to-family closure bridge. The local patch-deficit count gives the hard multiplicity $N_{\mathrm{fam}}=3$, the topological family step turns that multiplicity into a genuine family space, and the resulting occupancy $\Omega_{\mathrm{adm}}=48$ feeds back into the defect reading $\deltatop=\Omega_{\mathrm{adm}}\cthree^4$.}
\label{fig:family-triangle}
\end{figure}

The Hodge metric on $F$ turns the monodromy data into an honest unitary family system rather than a merely decorative triangle sketch.

\begin{theorem}[Topological family closure on the orbifold seam section]
\label{thm:topological-family-closure}
\label{thm:orbifold-aps-family-master}
Let $X_f$ be the orbifold seam section with four $\mathbb Z_4$ corners. Then
\[
\chi_{\mathrm{orb}}(X_f)=1-4\left(1-\frac14\right)=-2,
\qquad
\rank H^{\mathrm{free}}_1(X_f,\mathbb Z)=3.
\]
Hence
\[
F:=H^1(X_f,\mathbb C)\cong \mathbb C^3,
\qquad
N_{\mathrm{fam}}=3,
\qquad
\Oadm=N_{\mathrm{fam}}\dim S^+=3\cdot 16=48.
\]
Equivalently, the local patch-deficit count is
\[
N_{\mathrm{fam}}=F_{\mathrm{patch}}\!\left(1-\frac14\right)=4\cdot\frac34=3.
\]
\end{theorem}

\begin{remark}[How the family theorem is now used]
The family theorem is used at topological strength in the present version. Any orbifold APS
realization is read as compatible analytic structure on the same closure, not as an additional
main-text theorem burden.
\end{remark}

\begin{proposition}[Orbifold APS compatibility]
\label{prop:orbifold-aps-compatibility}
Under the additional analytical hypotheses of \Cref{app:aps},
\[
\Ind^{\mathrm{orb,APS}}(D_{\mathrm{rel}}^+\otimes L_F)=3.
\]
\end{proposition}

\begin{corollary}[Boundary-to-bulk bridge corollary for the admissible occupancy]
\label{cor:boundary-bulk-occupancy}
On the canonical branch,
\[
\dim(S^+\otimes F)=16\cdot 3=48=\Oadm.
\]
Hence the topological surcharge enters the present paper as a derived occupancy corollary with an
observation-facing reading:
\[
\deltatop
=
\cthree^4\,\dim(S^+\otimes F)
=
\Oadm\,\cthree^4.
\]
\end{corollary}

The family subsection is therefore no longer split between a local count, an independent index theorem, and a later occupancy reinterpretation. Those statements are compressed into one topological theorem with one bridge corollary, while the holonomy proposition below records the resulting unitary geometry on $F$.

This theorem is the family statement referenced later by the admissibility theorems: the closure layer uses it as the explicit family component of the master closure theorem rather than as an external import.

\begin{proposition}[Special unitary family holonomy on the determinant-trivial branch]
\label{prop:family-holonomy-su3}
Let $L_F\to X_f$ be the flat Hermitian family bundle with fiber
$F\cong H^1(X_f,\mathbb C)$. Assume that its determinant line is flat trivial,
\[
\det L_F \cong X_f\times \mathbb C,
\]
and that the chosen flat connection preserves this trivialization. Then
\[
\mathrm{Hol}(\nabla_F)\subset SU(3)_F.
\]
\end{proposition}

\begin{proof}
The Hermitian metric implies
\[
\mathrm{Hol}(\nabla_F)\subset U(3).
\]
The induced connection on $\det L_F$ has trivial holonomy by assumption, so for
every loop $\gamma$ one has
\[
\det \rho_F(\gamma)=1.
\]
Hence every holonomy matrix lies in $SU(3)$.
\end{proof}

The role of the present subsection is thus no longer only mnemonic. It displays both the local patch-deficit count and the global APS/Hodge realization before the later closure stack compresses them into the theorem-level compiler surface.

